\scanpage{45}
\begin{center}- 16 -\end{center}
\hfill n$^\circ$ 276

En effet, utilisons la filtration de $U(g)$ définie au n$^\circ$2, Exemple 3 ; elle est discrète : $U(g)$ est donc un anneau filtré séparé et complet ; or l'anneau gradué associé est un quotient d'un anneau de polynômes sur $K$, donc est noethérien ; d'où le corollaire.

7. - \uline{Complété d'un module filtré}.

Soit $E$ un groupe topologique abélien. Rappelons (Top. Gén.\unsure{.;}) que le complété $\hat E$ de l'espace uniforme $E$ est canoniquement muni d'une structure de groupe abélien, et que l'application canonique de $E$ dans $\hat E$ est une \uline{représentation} du groupe topologique $E$ dans le groupe topologique $\hat E$, dont le noyau est l'adhérence de 0 dans $E$.

Etant donné une \uline{représentation} $u$ de $E$ dans un groupe topologique abélien $E'$, l'application $u$ est uniformément continue, et par suite se prolonge en une application $\hat u:\hat E\longrightarrow\hat E'$ ; celle-ci est encore une \uline{représentation}

En particulier, soit $F$ un sous-groupe de $E$; l'application identique $j:F\longrightarrow E$ se prolonge en une représentation $\hat j$ du complété $\hat F$ dans le complété $\hat E$, laquelle est en fait un isomorphisme topologique de $\hat F$ sur un sous-groupe fermé de $\hat E$, à savoir l'adhérence de l'image de $F$ par l'application canonique de $E$ dans $\hat E$ ; nous identifierons toujours $\hat F$ à son image dans $\hat E$.

Considérons de plus le groupe topologique quotient $H=E/F$ ; la représentation canonique $E\xrightarrow{u}H$ se prolonge en une représentation $\hat E\xrightarrow{\hat u}H$, et comme $u$ applique $F$ sur $e$ il est clair que $\hat u$ applique $\hat F$ sur $e$ ; donc $\hat u$ définit par passage au quotient une représentation canonique
\[
v:\hat E/\hat F\longrightarrow\widehat{E/F}.
\]
On a d'autre part, à partir des représentations canoniques $E\to\hat E$ et $F\longrightarrow\hat F$, une représentation $E/F\to\hat E/\hat F$ et par suite une représentation
\mspage{44}
